\documentclass[11pt,a4paper]{article}
\usepackage[utf8]{inputenc}
\usepackage{amsmath}
\usepackage{amsfonts}
\usepackage{listings}
\usepackage{xcolor}
\usepackage{geometry}

% Page Geometry setup
\geometry{a4paper, margin=1in}

% Code snippet configurations
\definecolor{codegray}{rgb}{0.95,0.95,0.95}
\definecolor{commentgreen}{rgb}{0.12,0.53,0.22}
\lstset{
    backgroundcolor=\color{codegray},
    basicstyle=\ttfamily\small,
    breaklines=true,
    commentstyle=\color{commentgreen},
    frame=single,
    language=Python,
    showstringspaces=false
}

\title{\textbf{Analysis of Computational Graphs and Backpropagation}}
\author{}
\date{}

\begin{document}

\maketitle

To see exactly how this code works, let’s build a small mathematical graph, watch how the \texttt{backward()} function constructs it, and see how each line changes its state.

Let’s assume we are calculating a simple equation: 
\begin{equation*}
    z = (x + y) \times 2
\end{equation*}

\section{The Initial Computational Graph}
When you write expressions using your \texttt{Value} class, Python builds a directed graph behind the scenes. Every node stores its \textit{data}, its gradient (\textit{grad}), and pointer references to its creators (\texttt{\_prev}).

Here is what the initial state of our nodes looks like before \texttt{z.backward()} is called:
\begin{itemize}
    \item \textbf{Node $x$:} $\text{data}=3.0$, $\text{grad}=0.0$, $\texttt{\_prev}=()$
    \item \textbf{Node $y$:} $\text{data}=4.0$, $\text{grad}=0.0$, $\texttt{\_prev}=()$
    \item \textbf{Node $(x+y)$:} $\text{data}=7.0$, $\text{grad}=0.0$, $\texttt{\_prev}=(x, y)$
    \item \textbf{Node $z$:} $\text{data}=14.0$, $\text{grad}=0.0$, $\texttt{\_prev}=((x+y), 2)$
\end{itemize}

\section{How Each Section of \texttt{backward()} Updates the Graph}
Let's trace the function execution line-by-line to see how it dynamically alters this network.

\subsection*{Section A: The Setup}
\begin{lstlisting}
topo = []
visited = set()
\end{lstlisting}
\textbf{Effect on Graph:} No structural changes yet. This simply allocates a blank tracking sheet (\texttt{visited}) to avoid infinite loops, and an empty conveyor belt (\texttt{topo}) to structure our order of execution.

\subsection*{Section B: Building the Topological Order}
\begin{lstlisting}
def build(v):
    if v not in visited:
        visited.add(v)
        for child in v._prev:
            build(child)
        topo.append(v)

build(self) # triggered on 'z'
\end{lstlisting}
\textbf{Effect on Graph:} This walks \textit{backward} through the parent arrows (\texttt{\_prev}) to map out dependencies. Because it appends to \texttt{topo} after visiting children, it orders them from the very starting inputs up to the output.

\noindent \textbf{Resulting State of \texttt{topo}:} \texttt{[x, y, (x+y), Value(2), z]}

\subsection*{Section C: Planting the Seed Gradient}
\begin{lstlisting}
self.grad = 1.0
\end{lstlisting}
\textbf{Effect on Graph:} This updates the final node ($z$)'s gradient from $0.0$ to $1.0$. In calculus, the derivative of anything with respect to itself is 1:
\begin{equation*}
    \frac{\partial z}{\partial z} = 1
\end{equation*}
Without this line, all downstream gradients would multiply by zero and vanish.

\subsection*{Section D: The Domino Effect (Backpropagation)}
\begin{lstlisting}
for node in reversed(topo):
    node._backward()
\end{lstlisting}
By reversing our \texttt{topo} list, we loop through the nodes in exact reverse order: \texttt{[z, Value(2), (x+y), y, x]}.

As each node calls its specific \texttt{.\_backward()} function, it applies the mathematical \textbf{Chain Rule} to update the gradients of its immediate parent inputs:

\subsubsection*{Processing $z$:}
\begin{itemize}
    \item It looks at its operation ($\times$).
    \item It updates its inputs' gradients: \texttt{(x+y).grad} becomes $2.0$ (since $1.0 \times 2$).
\end{itemize}

\subsubsection*{Processing $(x+y)$:}
\begin{itemize}
    \item It looks at its operation ($+$).
    \item It passes its accumulated gradient ($2.0$) directly back to its inputs: \texttt{x.grad} becomes $2.0$ and \texttt{y.grad} becomes $2.0$.
\end{itemize}

\section*{Final State of the Graph}
Once the loop finishes executing, your computational graph's state is completely transformed. The forward data remains untouched, but the analytical derivative values have cleanly saturated throughout the entire architecture:
\begin{itemize}
    \item \textbf{Node $x$:} $\text{data}=3.0$, $\text{grad}=2.0$
    \item \textbf{Node $y$:} $\text{data}=4.0$, $\text{grad}=2.0$
    \item \textbf{Node $(x+y)$:} $\text{data}=7.0$, $\text{grad}=2.0$
    \item \textbf{Node $z$:} $\text{data}=14.0$, $\text{grad}=1.0$
\end{itemize}

\end{document}